% Section 13: benchmark results and failure analysis.
\section{Benchmark results}\label{sec:results}

\pgfplotsset{discard if not/.style 2 args={
  x filter/.append code={\edef\tempa{\thisrow{#1}}\edef\tempb{#2}\ifx\tempa\tempb\else\def\pgfmathresult{inf}\fi}}}

Three public test sets are used, all with published optimal values: the \nSpxTot{} LPs of
Netlib~\cite{gay1985}, from $\netlibMinRows$ to $\netlibMaxRows$ rows and up to \netlibMaxNnz{}
nonzeros (\code{\netlibMaxName}); the \nMipList{} MILPs of MIPLIB~3~\cite{bixby1998}; and the
\nQpList{} convex QPs of Maros and Mészáros~\cite{maros1999}. The large LPs of \cref{sec:pdlp} come
from Mittelmann's benchmark sets~\cite{mittelmann}. Every number below is read from the result files
named in the captions. \Cref{tab:summary} gives the overview and \cref{fig:bars} the solved counts.

\ifnum\haveInProgress=1
\begin{keybox}[Sweeps still running when this report was built]
\file{gen\_data.py} analyses the newest sweep of each suite that covers every instance. The following
newer sweeps were still running and are not yet in the headline numbers: \inProgressList. Where a
comparison with them is meaningful on the instances finished so far, the text gives it and says so.
Rerunning \file{gen\_data.py} and \code{latexmk} after they finish updates every number.
\end{keybox}
\fi

\begin{table}[tbp]
\centering\small
\caption[Summary of the benchmark sweeps]{Summary of the benchmark sweeps. ``Solved'' uses the
criteria of \cref{sec:verification}. Times are shifted geometric means (shift \shiftA\,s) over the
instances \nirnay{} solves, for \nirnay{} and for the reference on the same instances.}
\label{tab:summary}
\begin{tabular}{@{}l l l r r r r@{}}
\toprule
Suite & Engine & Reference & Solved & \nirnay{} (s) & Ref.\ (s) & Ratio\\
\midrule
Netlib LP & dual simplex (\spxVer) & HiGHS & \nSpx/\nSpxTot & \sgmSpx & \sgmSpxRef & \ratioSpx\\
Netlib LP & interior point (\ipmVer) & HiGHS & \nIpm/\nIpmTot & \sgmIpm & \sgmIpmRef & \ratioIpm\\
Netlib LP & PDLP, $10^{-4}$ & HiGHS & \nPdlp/\nPdlpTot$^{a}$ & \sgmPdlp & \sgmPdlpRef & \ratioPdlp\\
MIPLIB 3 & branch-and-bound (\mipVer) & HiGHS & \nMip/\nMipTot$^{b}$ & \sgmMip & \sgmMipRef & \ratioMip\\
Maros--Mészáros & QP interior point & Clarabel & \nQp/\nQpTot & \sgmQp & \sgmQpRef & \ratioQp\\
\bottomrule
\multicolumn{7}{@{}l}{\footnotesize $^a$ status \code{optimal} at relative KKT $10^{-4}$; objective errors in \cref{sec:pdlp-results}.\quad
$^b$ \tlMip\,s limit; HiGHS proves \nMipRefOpt.}
\end{tabular}
\end{table}

\begin{figure}[tbp]
\centering
\begin{tikzpicture}
\begin{axis}[
  width=0.82\linewidth, height=5.2cm,
  xbar, bar shift=0pt, bar width=9pt,
  xmin=0, enlarge y limits=0.14,
  ytick=data, yticklabels from table={data/solved_bars.dat}{label},
  y dir=reverse,
  xlabel={instances},
  xmajorgrids,
  nodes near coords, nodes near coords style={font=\sffamily\scriptsize},
  every node near coord/.append style={anchor=west},
  legend style={at={(1,1.03)},anchor=south east}, legend columns=2,
]
\addplot[fill=rule,draw=muted] table[x=total,y=x] {data/solved_bars.dat};
\addplot[fill=saffron,draw=saffron!80!black,nodes near coords style={text=white,anchor=east}] table[x=solved,y=x] {data/solved_bars.dat};
\legend{instances run, solved}
\end{axis}
\end{tikzpicture}
\caption[Solved instances per suite]{Solved instances per suite and engine (from
\file{data/solved\_bars.dat}, generated from \file{results/}). PDLP counts status \code{optimal} at its
$10^{-4}$ tolerance.}
\label{fig:bars}
\end{figure}

\subsection{Netlib LP: correctness}
The dual simplex solves all \nSpx{} of the \nSpxTot{} Netlib problems to the reference objective
(file \file{\spxFile}). The median relative objective error is \medErrSpx{} and the largest \maxErrSpx{},
so the answers are vertex solutions exact to rounding. The interior-point method solves \nIpm{} (median
error \medErrIpm, largest \maxErrIpm). \Cref{fig:ecdf} shows the distribution of the errors for every
engine. The interior-point failures are analysed in \cref{sec:ipm-fail}.

\begin{figure}[tbp]
\centering
\begin{tikzpicture}
\begin{axis}[
  width=0.9\linewidth, height=6cm,
  xlabel={$\log_{10}$ relative objective error against the reference},
  ylabel={fraction of scored instances},
  xmin=-17.2, xmax=0, ymin=0, ymax=1.02,
  grid=major, legend cell align=left,
  legend style={at={(0.5,-0.22)},anchor=north}, legend columns=2,
]
\addplot[nspx,const plot] table[x=logerr,y=frac] {data/ecdf_spx.dat};
\addplot[nipm,const plot] table[x=logerr,y=frac] {data/ecdf_ipm.dat};
\addplot[npdlp,const plot] table[x=logerr,y=frac] {data/ecdf_pdlp.dat};
\addplot[igreen,line width=1.1pt,const plot] table[x=logerr,y=frac] {data/ecdf_qp.dat};
\addplot[navy,line width=1.1pt,densely dashed,const plot] table[x=logerr,y=frac] {data/ecdf_mip.dat};
\addplot[muted,dashed,forget plot] coordinates {(-6,0) (-6,1.02)};
\addplot[muted,dashed,forget plot] coordinates {(-4,0) (-4,1.02)};
\legend{dual simplex (Netlib), interior point (Netlib), {PDLP, status optimal (Netlib)}, QP interior point (Maros--Mészáros), branch-and-bound (MIPLIB 3)}
\end{axis}
\end{tikzpicture}
\caption[Distribution of relative objective errors]{Cumulative distribution of the relative
objective error \eqref{eq:relerr} over the instances each engine reports optimal (for LP, QP and the
dual simplex: the instances counted as solved; for PDLP: every status \code{optimal}). Errors of
exactly zero are plotted at $10^{-17}$. Dashed lines: the $10^{-6}$ LP/QP and $10^{-4}$ MILP criteria.}
\label{fig:ecdf}
\end{figure}

\subsection{Netlib LP: speed}
\Cref{fig:profile} shows the performance profiles of the dual simplex, the interior-point method and
HiGHS on Netlib, and \cref{fig:scatter} plots every instance's time against the reference's.

\begin{honestbox}[Speed, stated plainly]
\nirnay{}'s dual simplex is faster than HiGHS on \nFasterSpx{} of the \nSpx{} Netlib problems it solves.
Its shifted geometric mean time is \sgmSpx\,s against \sgmSpxRef\,s, a ratio of \ratioSpx{} with a
\shiftA\,s shift and \ratioSpxTenShift{} with a \shiftB\,s shift. The per-instance ratios have
geometric mean \gmrSpx{}, median \medrSpx{} and maximum \maxrSpx{}. The shifted mean understates the
gap because most Netlib problems solve in milliseconds, so the per-instance ratios are the fairer
summary: \nirnay{} is about \gmrSpx{} times slower. On load-normalised times (\cref{fig:profile}), the
interior-point method is faster than the simplex on \nIpmFasterThanSpx{} of the \nBothSolved{} problems
both solve, and taking the faster of the two engines on each problem gives a shifted-mean ratio of
\bestOfTwoRatio{} against HiGHS.
\end{honestbox}

\begin{figure}[tbp]
\centering
\begin{tikzpicture}
\begin{axis}[
  width=0.9\linewidth, height=6.2cm,
  xmode=log, log basis x=2,
  xmin=1, xmax=\profTauMax, ymin=0, ymax=1.02,
  xlabel={$\tau$ (time relative to the fastest of the solvers shown)},
  ylabel={$\rho_s(\tau)$},
  grid=major, legend pos=south east, legend cell align=left,
  log number format basis/.code 2 args={\pgfmathparse{#1^#2}\pgfmathprintnumber{\pgfmathresult}},
]
\addplot[nhighs,const plot] table[x=tau,y=rho] {data/profile_highs.dat};
\addplot[nspx,const plot] table[x=tau,y=rho] {data/profile_spx.dat};
\addplot[nipm,const plot] table[x=tau,y=rho] {data/profile_ipm.dat};
\legend{HiGHS, \nirnay{} dual simplex, \nirnay{} interior point}
\end{axis}
\end{tikzpicture}
\caption[Performance profile on Netlib]{Dolan--Moré performance profile on the \nProf{} Netlib
problems. Times come from \file{\spxFile} and \file{\ipmFile}. The two sweeps ran under different
machine load (HiGHS itself was \sweepLoadRatio{} times slower, in geometric mean, in the first than in
the second), so each engine's time is rescaled by the ratio of the common HiGHS time (the geometric
mean of its two measurements) to HiGHS's time in that engine's sweep. HiGHS is
fastest on a fraction \rhoOneHighs{} of the problems. The dual simplex is within a factor of 10 of the
fastest on a fraction \rhoTenSpx{} and within \tauAllSpx{} on all of them. The interior-point
method's curve stays below~1 because of its failures.}
\label{fig:profile}
\end{figure}

\begin{figure}[tbp]
\centering
\begin{tikzpicture}
\begin{groupplot}[
  group style={group size=2 by 2, horizontal sep=14mm, vertical sep=16mm},
  width=0.49\linewidth, height=0.47\linewidth,
  xmode=log, ymode=log,
  xlabel={reference time (s)}, ylabel={\nirnay{} time (s)},
  grid=major, unbounded coords=discard,
  title style={font=\sffamily\footnotesize\bfseries,color=navy},
]
\nextgroupplot[title={Netlib: dual simplex vs.\ HiGHS}, xmin=1e-4, xmax=1e3, ymin=1e-4, ymax=1e3]
  \addplot[muted,dashed,domain=1e-4:1e3,samples=2,forget plot] {x};
  \addplot[muted!60,dotted,domain=1e-4:1e2,samples=2,forget plot] {10*x};
  \addplot[only marks,mark=*,mark size=1.3pt,color=saffron,discard if not={ok}{1}] table[x=ref,y=time] {data/scatter_spx.dat};
  \addplot[only marks,mark=x,mark size=2.4pt,color=navy,thick,discard if not={ok}{0}] table[x=ref,y=time] {data/scatter_spx.dat};
\nextgroupplot[title={Netlib: interior point vs.\ HiGHS}, xmin=1e-4, xmax=1e3, ymin=1e-4, ymax=1e3]
  \addplot[muted,dashed,domain=1e-4:1e3,samples=2,forget plot] {x};
  \addplot[muted!60,dotted,domain=1e-4:1e2,samples=2,forget plot] {10*x};
  \addplot[only marks,mark=*,mark size=1.3pt,color=sky,discard if not={ok}{1}] table[x=ref,y=time] {data/scatter_ipm.dat};
  \addplot[only marks,mark=x,mark size=2.4pt,color=navy,thick,discard if not={ok}{0}] table[x=ref,y=time] {data/scatter_ipm.dat};
\nextgroupplot[title={MIPLIB 3: branch-and-bound vs.\ HiGHS}, xmin=1e-2, xmax=2e2, ymin=1e-2, ymax=2e2]
  \addplot[muted,dashed,domain=1e-2:2e2,samples=2,forget plot] {x};
  \addplot[muted!60,dotted,domain=1e-2:2e1,samples=2,forget plot] {10*x};
  \addplot[only marks,mark=*,mark size=1.3pt,color=igreen,discard if not={ok}{1}] table[x=ref,y=time] {data/scatter_mip.dat};
  \addplot[only marks,mark=x,mark size=2.4pt,color=navy,thick,discard if not={ok}{0}] table[x=ref,y=time] {data/scatter_mip.dat};
\nextgroupplot[title={Maros--Mészáros: QP-IPM vs.\ Clarabel}, xmin=1e-4, xmax=1e3, ymin=1e-4, ymax=1e3]
  \addplot[muted,dashed,domain=1e-4:1e3,samples=2,forget plot] {x};
  \addplot[muted!60,dotted,domain=1e-4:1e2,samples=2,forget plot] {10*x};
  \addplot[only marks,mark=*,mark size=1.3pt,color=plum,discard if not={ok}{1}] table[x=ref,y=time] {data/scatter_qp.dat};
  \addplot[only marks,mark=x,mark size=2.4pt,color=navy,thick,discard if not={ok}{0}] table[x=ref,y=time] {data/scatter_qp.dat};
\end{groupplot}
\end{tikzpicture}
\caption[Time against reference time]{Time of \nirnay{} against the time of the reference solver, per
instance, on log--log axes. Dots: solved; crosses: not solved (plotted at the time used). Dashed:
equal time; dotted: ten times slower. Points above the dashed line are instances where \nirnay{} is
slower. On MIPLIB the \tlMip\,s limit shows as a horizontal row of crosses.}
\label{fig:scatter}
\end{figure}

\paragraph{Where the time goes}
The iteration counts are reasonable: the median Netlib solve takes \medItSpx{} dual iterations, and
\cref{fig:iters} shows the count growing roughly in proportion to $m+n$, as expected for a dual simplex.
The cost is in each iteration. Every step calls several compiled kernels (pricing, BTRAN, pivot row,
ratio test, two or three FTRANs, updates) from a Python loop. On small problems the Python overhead
of each call exceeds the arithmetic. On large ones the linear algebra is less refined than in
production codes: FTRAN and BTRAN work on dense vectors of length $m$ instead of exploiting
hyper-sparsity, the product-form update grows the eta file where a Forrest--Tomlin update would
modify $U$ in place. \ifnum\spxHasPresolve=1 The simplex sweep analysed here runs with presolve,
which is still written in plain Python; its cost and benefit are measured in \cref{sec:presolve}.\else
The sweeps analysed here ran without presolve; its effect is measured in \cref{sec:presolve}.\fi{} \Cref{sec:roadmap} lists the remedies in the order we expect them to pay.

\begin{figure}[tbp]
\centering
\begin{tikzpicture}
\begin{axis}[
  width=0.62\linewidth, height=5.2cm, xmode=log, ymode=log,
  xlabel={$m+n$ (rows + columns)}, ylabel={dual simplex iterations},
  grid=major, legend pos=north west, legend cell align=left,
]
\addplot[only marks,mark=*,mark size=1.2pt,color=saffron] table[x=mn,y=iters] {data/iters_spx.dat};
\addplot[muted,dashed,domain=40:30000,samples=2] {x};
\legend{Netlib problems, iterations $=m+n$}
\end{axis}
\end{tikzpicture}
\caption[Dual simplex iterations against problem size]{Dual simplex iterations (phase~1, phase~2
and cleanup together) against $m+n$ on the Netlib problems.}
\label{fig:iters}
\end{figure}

\subsection{Interior-point failures}\label{sec:ipm-fail}
\begin{table}[ht]
\centering\small
\caption[Interior-point failures on Netlib]{Netlib problems the interior-point method does not solve (\file{\ipmFile}).}
\label{tab:ipm-fail}
\begin{tabular}{@{}l rr l r r r@{}}
\toprule
Problem & $m$ & $n$ & status & rel.\ err & time (s) & iterations\\
\midrule
\rows{data/ipm_failures.tex}
\bottomrule
\end{tabular}
\end{table}
\Cref{tab:ipm-fail} lists the failures. All end in the stall rule's
\code{numerical\_error}, and \nIpmNear{} of them are within $10^{-3}$ of the optimum when they stop. These
are the problems known to be hard for interior-point methods without presolve. The comparator study
found Clarabel and GLPK's interior point also failing on \code{greenbea} without presolve. The
\code{pilot} models have coefficients spanning many orders of magnitude, and \code{dfl001}'s normal
matrix fills heavily. Production IPMs solve them with presolve, crossover to a basis, and
more careful handling of free variables. Within \nirnay{}, the dual simplex solves all of them.

\subsection{MIPLIB 3}\label{sec:mip-results}
\ifnum\nMipTot<\nMipList
\begin{keybox}[Partial sweep]
The MIPLIB sweep (\file{\mipFile}) had finished \nMipTot{} of the \nMipList{} instances when this
report was built. The numbers below cover those instances and update when \file{gen\_data.py} is rerun.
\end{keybox}
\fi
With a \tlMip\,s limit, branch-and-bound proves \nMip{} of \nMipTot{} instances optimal, all with the
reference objective. No instance was declared optimal with a wrong objective (\nMipWrongOpt{}). HiGHS,
under the same limit, proves \nMipRefOpt. \nMipBothOpt{} instances are proved by both, \nMipRefOnly{}
only by HiGHS, and \nMipOnlyUs{} only by \nirnay{}. \nirnay{} is faster than HiGHS on \nFasterMip{}
of its solved instances (\fasterMipNames). On the instances it solves, its shifted geometric mean time is
\sgmMip\,s against \sgmMipRef\,s for HiGHS, \ratioMip{} times slower.

The instances not proved fall into three groups:
\begin{itemize}
  \item \textbf{optimum found, not proved} (\nMipFoundNotProved): \mipFoundNotProvedNames. The incumbent is
        optimal, but the bound does not close in time;
  \item \textbf{incumbent with a gap} (\nMipIncumbent{} with an incumbent at the limit, of which
        \nMipGapBelowOnePct{} are within $1\%$ and \nMipGapBelowTenPct{} within $10\%$), see \cref{fig:gaps};
  \item \textbf{no incumbent at all} (\nMipNoIncumbent): \mipNoIncumbentNames. These are mostly
        large set-partitioning and network-design models, where rounding and diving fail and each LP is
        expensive.
\end{itemize}
The causes are the ones the comparator study (\cref{sec:competitors}) points to. \nirnay{}'s MIP presolve
only tightens bounds. Its only cut family is Gomory mixed-integer; production codes add MIR, knapsack
cover, flow cover, clique and zero-half cuts. It has no feasibility pump, no feasibility jump and no
large-neighbourhood search. Node throughput is low: the median solved instance processes
\mipNodesPerSec{} nodes per second at a median of \mipLpItersPerNode{} dual iterations per node, and
each node pays the Python overhead of the simplex driver. The largest search on record took
\maxNodesMip{} nodes.

\begin{figure}[tbp]
\centering
\begin{tikzpicture}
\begin{axis}[
  width=0.8\linewidth, y=10.5pt,
  xbar, bar width=6pt, xmode=log, log origin=infty,
  xmin=1e-4, xmax=2,
  ytick=data, yticklabels from table={data/mip_gaps.dat}{name},
  yticklabel style={font=\ttfamily\scriptsize},
  xlabel={relative gap at the time limit}, xmajorgrids, enlarge y limits={abs=0.7},
  y dir=reverse,
]
\addplot[fill=saffron,draw=saffron!80!black] table[x=gap,y=idx] {data/mip_gaps.dat};
\end{axis}
\end{tikzpicture}
\caption[MIPLIB 3: gap at the time limit]{MIPLIB~3 instances that reached the \tlMip\,s limit with an
incumbent: relative gap $|f-\underline z|/\max(1,|f|)$ at the limit (log scale). The target is $10^{-4}$.}
\label{fig:gaps}
\end{figure}

\ifnum\haveMipVtwo=1
\paragraph{Effect of cuts and presolve} The sweep with Gomory cuts and presolve
(\file{\mipNewerFile})\ifnum\mipNewerPartial=1{} was still running when this report was built, with
\mipNewerDone{} of \nMiplibFiles{} instances finished\fi. It proves \mipNewerSolved{} of them optimal.
On the \nMipCommon{} instances both sweeps cover, it proves \nMipCutsGained{} instances that v1 (no
cuts, no presolve) did not (\mipGainedNames), and loses \nMipCutsLost{} (\mipLostNames).
\fi

\subsection{Maros--Mészáros QP}
The QP interior point solves \nQp{} of the \nQpTot{} problems (\file{\qpFile}), with a median error of
\medErrQp{} and a median of \medItQp{} iterations (the most \maxItQp). It is faster than Clarabel on
\nFasterQp{} of them, and slower by a factor of \ratioQp{} in shifted geometric mean. The rest divide as
follows:
\begin{itemize}
  \item \textbf{reference failed} (\nQpRefBad): \qpRefBadNames. Clarabel did not report a solution,
        so these count as not solved whatever \nirnay{} returned;
  \item \textbf{accepted loosely} (\nQpLoose): \qpLooseNames. \nirnay{} reports \code{optimal} with an
        objective error between $10^{-6}$ and $10^{-3}$ against the reference: points whose relative
        residuals met the tolerance while the objective error did not (or, see below, reference error);
  \item \textbf{wrong answer} (\nQpWrongBig): \qpWrongBigNames. Reported \code{optimal} with a large
        objective error against the reference;
  \item \textbf{numerical error} (\nQpNumErr): \qpNumErrNames. The stall rule gave up. Clarabel itself
        reports only \code{AlmostSolved} on \nQpAlmostRef{} of the problems \nirnay{} does not solve;
  \item \textbf{other} (\nQpOther): \qpOtherNames. The solve stopped at the time limit, or the
        benchmark process returned no status; these count as not solved.
\end{itemize}
\ifnum\nQpLooseOrWrong>0 The wrong and loose answers show the weak point of the QP method: its
termination rule trusts relative residuals that do not bound the objective error on badly scaled
problems.\else No problem is reported \code{optimal} with an objective error above $10^{-6}$: every
failure is an honest status (numerical error, time limit) rather than a wrong answer.\fi{}
\ifnum\nQpRefAnnotated>0 On \nQpRefAnnotated{} files the two MPS parsers disagree, and the reference
solves \nirnay{}'s reading (\cref{sec:verification}).\fi{}
\ifnum\qpVerNum=1 Not every entry above is \nirnay{}'s fault. This sweep used the first version of the
QP reference, whose two defects are described in \cref{sec:verification}: Clarabel at default
tolerances, and HiGHS's parser misreading \code{DPKLO1} and dropping small coefficients in
\code{HUESTIS}, \code{HUES-MOD} and \code{KSIP}. Some ``loose'' and ``wrong'' entries, \code{DPKLO1} among
them, may be reference errors, and the corrected reference will tell. The changes described in \cref{sec:qp} (the gap measured with the objective
constant, the residual check on every direction, and the floored starting point) were made in response
to this sweep. \ifnum\haveQpNew=1 The sweep that measures them was still running when this report was
built.\fi\fi{} Presolve is not yet applied to QPs. \Cref{tab:maros-full} in the appendix lists every
problem.
